\documentclass{article}
\usepackage{amsmath, amssymb, amsthm}
\usepackage{hyperref}
\usepackage{geometry}
\geometry{margin=1in}

\title{Erd\H{o}s Problem \#1100}
\date{Last updated: 2025-10-19}
\author{Source: erdosproblems.com}

\begin{document}
\maketitle

\noindent\textbf{Status:} OPEN \quad \textbf{No prize}

\noindent\textbf{Tags:} number theory, divisors

\medskip

\noindent\textbf{Problem Statement:}

If $1=d_1&#60;\cdots&#60;d_{\tau(n)}=n$ are the divisors of $n$, then let $\tau_\perp(n)$ count the number of $i$ for which $(d_i,d_{i+1})=1$. Is it true that $\tau_\perp(n)/\omega(n)\to \infty$ for almost all $n$? Is it true that\[\tau_\perp(n)&#60; \exp((\log n)^{o(1)})\]for all $n$? Let\[g(k) = \max_{\omega(n)=k}\tau_\perp(n),\]where $\omega(n)$ counts the number of distinct prime divisors of $n$, and $n$ is restricted to squarefree integers. Determine the growth of $g(k)$.




The function $\tau_\perp(n)$ was considered by Erd\H{o}s and Hall \cite{ErHa78}. It is trivial that $\tau_\perp(n)\geq \omega(n)$ (with equality for infinitely many $n$). Erd\H{o}s and Hall prove, for all $\epsilon&#62;0$ and sufficiently large $x$,\[\max_{n&lt;x} \tau_\perp(n) &gt; \exp((\log\log x)^{2-\epsilon}).\]Erd\H{o}s and Simonovits (see \cite{Er81h}) proved\[(2^{1/2}+o(1))^k &lt; g(k) &lt; (2-c)^k\]for some constant $c&gt;0$.




References

[Er81h] Erd\H{o}s, P., Some problems and results on additive and multiplicative number theory . Analytic number theory (Philadelphia, Pa., 1980) (1981), 171-182.

[ErHa78] Erd\H{o}s, P. and Hall, R. R., On some unconventional problems on the divisors of integers . J. Austral. Math. Soc. Ser. A (1978), 479--485.

\noindent\textbf{Related OEIS sequences:} A325864, possible


\bigskip
\noindent\small{Source: \url{https://www.erdosproblems.com/1100}}
\end{document}
